\begin{textsample}{POS Dim 4 – 2017\_c6 – Score 11.00 – 2017\_c6/enade\_2017\_admin\_1e32d435fe.txt}  \label{ex:f4_pos_034}
CAPÍTULO 6 ANÁLISE TÉCNICA DA PROVA Este capítulo tem por objetivo apresentar o desempenho dos estudantes concluintes de Arquitetura e Urbanismo no Enade/2017.
Para isso, foram calculadas as estatísticas básicas da prova em seu todo ( seção 6.1.1 ), bem como as estatísticas relacionadas aos Componentes de Formação Geral ( seção 6.1.2 ) e de Conhecimento Específico da Área ( seção 6.1.3 ).
Dadas as suas características, foram analisadas, em separado, as questões objetivas ( seção 6.2 ) e as questões discursivas ( seção 6.3 ).
Para as questões objetivas, foram disponibilizados os índices de facilidade e de discriminação ponto-bisserial, também em separado, para os Componentes de Formação Geral ( seção 6.2.1 ) e de Conhecimento Específico ( seção 6.2.2 ).
De cada componente, uma das questões foi escolhida para exemplificar a análise gráfica, relacionando as alternativas escolhidas pelos estudantes ( inclusive o gabarito ) com o número de acertos no componente.
O Anexo I apresenta a íntegra da análise gráfica para todas as questões objetivas.
Para cada uma das questões discursivas, os conteúdos dos tipos mais comuns de respostas dos estudantes são apresentados e comparados com o padrão de respostas esperado ( ver Anexo VIII com o padrão de respostas ).
Tomando como base as duas questões discursivas do Componente de Formação Geral, a seção 6.3.1.6 apresenta comentários sobre a correção das respostas com respeito à Língua Portuguesa.
As tabelas apresentam as seguintes estatísticas das notas19 : média do desempenho na prova, erro padrão da média, desvio padrão, nota mínima, mediana e nota máxima para cada um de seus componentes.
Tais estatísticas contemplam o total de estudantes concluintes da Área de Arquitetura e Urbanismo inscritos e presentes à prova do Enade/2017, tendo em vista agregações, ou por Grandes Regiões e o país como um todo, ou por Categoria Administrativa, Organização Acadêmica da IES e Modalidade de Ensino.
Em relação aos gráficos de distribuição de notas, o intervalo considerado foi de 10 unidades, aberto à esquerda e fechado à direita, com exceção do primeiro intervalo, [ 0 ; 10 ], fechado em ambos os extremos.
Para os gráficos de distribuição das notas das questões discursivas, foram consideradas mais duas categorias : questão em branco20 e nota zero. 19 Uma definição dessas estatísticas pode ser encontrada no Glossário. 20 Nesse grupo estão incluídas também as respostas classificadas como nulas ou desconsideradas. 6.1 ESTATÍSTICAS BÁSICAS DA PROVA Esta seção apresenta estatísticas selecionadas e histogramas da nota geral ( 6.1.1 ) e de cada componente : Formação Geral ( 6.1.2 ) e Conhecimento Específico ( 6.1.3 ).
São também apresentadas estatísticas selecionadas de subpopulações caracterizadas por Grande Região, Categoria Administrativa, Organização Acadêmica e Modalidade de Ensino ( quando pertinente ). 6.1.1 Estatísticas Básicas Gerais A Tabela 6.1 apresenta as Estatísticas Básicas da prova, por Grande Região, dos estudantes concluintes de Arquitetura e Urbanismo.
A Média das notas da prova como um todo ( nas seções seguintes serão analisados os Componentes de Formação Geral e de Conhecimento Específico ) foi 49,0, sendo que os alunos da região Sul obtiveram a Média mais alta ( 51,4 ), e os da região Norte obtiveram a Média mais baixa ( 43,0 ).
As demais médias foram : 49,4 na região Nordeste, 49,0 na região Sudeste e 47,2 na região Centro-Oeste.
O Desvio padrão para o Brasil como um todo foi 14,4, sendo o maior Desvio padrão encontrado na região Nordeste ( 15,1 ) e o menor, nas regiões Norte ( 13,4 ), indicando uma dispersão um pouco menor das notas desta última região.
A região que obteve a maior nota Máxima foi a Sul ( 92,8 ), ao passo que a região que atingiu a menor nota Máxima foi a Norte ( 83,3 ).
A Mediana do Brasil como um todo foi 49,1, sendo a maior Mediana obtida na região Sul ( 51,5 ), e a menor obtida na Norte ( 42,4 ).
A nota Mínima foi zero em todas as regiões, com exceção da região Norte, onde a nota Mínima foi 9,3.
Considerando -se as notas segundo Grande Região, observa -se que existe diferença estatisticamente significativa ao nível de 95 \% entre a menor Média, obtida na região Norte ( 43,0 ), e as médias de todas as demais regiões21. 21 Todas as comparações deste capítulo utilizam os intervalos de 95 \%.
Os erros-padrão da média que possibilitam os testes estão disponíveis nas tabelas.
O comportamento das notas dos estudantes de todo o Brasil pode ser observado no Gráfico 6.1 que apresenta um histograma com a distribuição das mesmas.
Essa é uma distribuição unimodal com a moda no intervalo ( 40 ; 50 ], seguido pelo intervalo ( 50 ; 60 ].
A Tabela 6.2 apresenta informações referentes à Média da nota final desagregadas por Categoria Administrativa e por Organização Acadêmica.
A Média das notas da prova como um todo foi 49,0.
Em relação à Categoria Administrativa, os estudantes das IES Públicas obtiveram Média mais alta ( 58,9 ), e os das IES Privadas obtiveram Média mais baixa ( 46,9 ), que a Média nacional.
Observa -se que existe diferença estatisticamente significativa entre as médias das notas das IES Públicas e Privadas.
A diferença entre as médias das regiões Sul e Norte ( 8,4 ), a maior e a menor Média, é inferior à diferença entre IES Públicas e Privadas ( 12,0 ), caracterizando uma maior diversidade administrativa do que regional.
O Desvio padrão para as IES Públicas ( 14,1 ) foi inferior ao do Brasil como um todo ( 14,4 ), indicando uma dispersão um pouco menor das notas nesta Categoria Administrativa.
No tocante à Organização Acadêmica, as Faculdades e os Centros Universitários obtiveram Média mais baixa que a nacional ( 46,2 e 44,5, respectivamente ).
A Média das Universidades e dos CEFET/IFET foram maiores, respectivamente 51,1 e 53,3.
Constata -se que existe diferença estatisticamente significativa ao nível de 95 \% entre os valores mais altos da Média, o das Universidades e o dos CEFET/IFET, e os dos outros dois tipos de Organização Acadêmica, com valores mais baixos.
Não foi oferecido nenhum cursos na modalidade a Distância na área de Arquitetura e Urbanismo.
A Média da Modalidade Educação Presencial ( 49,0 ) foi, portanto, igual à Média do total de cursos. 6.1.2 Estatísticas Básicas no Componente de Formação Geral A Tabela 6.3 apresenta as Estatísticas Básicas em relação ao componente da prova que avalia a Formação Geral dos estudantes concluintes.
Os alunos de todo o Brasil obtiveram desempenho médio de 54,1.
Quanto à variabilidade, o Desvio padrão das notas dos estudantes do Brasil, como um todo, foi 17,2.
A maior Média foi obtida na região Sul ( 56,1 ), e a menor, na região Norte ( 49,6 ).
As demais médias foram : 54,3 na região Nordeste, 54,1 na região Sudeste e 53,0 na A maior nota no Componente de Formação Geral da prova do Enade/2017 foi 98,8, obtida por, pelo menos, um aluno nas regiões Sudeste e Sul.
A menor nota Máxima foi obtida na região Norte ( 94,6 ).
A Mediana do Brasil, como um todo, foi 52,2, sendo a menor Mediana encontrada na região Norte ( 50,8 ), e a maior, encontrada na região Sul ( 57,3 ).
A nota Mínima nesta parte foi zero em todas as regiões, sem exceção.
Considerando -se as notas segundo Grande Região, observa -se que existe diferença estatisticamente significativa entre a maior Média das notas do Componente de Formação Geral, obtida na região Sul ( 56,1 ), e todas as demais.
Também foram estatisticamente significativas as diferenças entre a menor Média, a da região Norte, e as demais.
O Gráfico 6.2 propicia a avaliação do desempenho dos estudantes do Componente de Formação Geral, a partir do histograma da distribuição das notas correspondentes.
A distribuição é unimodal, com moda em ( 50 ; 60 ], acima do intervalo modal da distribuição de notas da prova, como um todo ( Gráfico 6.1 ).
Nota -se, ainda que no Gráfico 6.2, as notas apresentam maior dispersão quando comparada com a do Gráfico 6.1 ( distribuição das notas da prova ), confirmada pela comparação dos desvios padrões : 14,4 para a nota da prova como um todo e 17,2 para o Componente de Formação Geral.
Na Tabela 6.4 são apresentadas as informações referentes ao desempenho dos concluintes do Componente de Formação Geral, em diferentes agregações : Categoria Administrativa, Organização Acadêmica e Modalidade de Ensino.
Observa -se que existe diferença estatisticamente significativa entre as médias dos tipos de Categoria Administrativa : a maior Média obtida por estudantes de IES Públicas ( 61,9 ) e a menor pelos de IES Privadas ( 52,5 ).
Nota -se que também há diferença estatisticamente significativa entre as médias dos quatro tipos de Organização Acadêmica, valor maior para os CEFET/IFET ( 59,8 ), seguidos pelas Universidades ( 55,9 ), pelos Centros Universitários ( 51,9 ) e pelas Faculdades ( 50,4 ). 6.1.3 Estatísticas Básicas do Componente de Conhecimento Específico A Tabela 6.5 apresenta as Estatísticas Básicas referentes ao Componente de Conhecimento Específico da área de Arquitetura e Urbanismo.
A Média do desempenho dos alunos do Brasil, como um todo, foi 47,3.
A maior Média foi obtida na região Sul ( 49,9 ), e a menor, na região Norte ( 40,8 ).
As demais médias foram : 47,7 na região Nordeste, 47,2 na região Sudeste e 45,2 na região Centro-Oeste.
Quanto à variabilidade das notas, o Desvio padrão do Brasil, como um todo, foi 15,8, sendo o maior Desvio padrão observado nas regiões Nordeste e Centro-Oeste ( 16,3 ), e o menor, na região Norte ( 14,6 ).
Os demais desvios foram : 15,4 na região Sudeste e 15,5 na Sul.
A Mediana das notas dos estudantes de todo o Brasil foi 47,4.
A maior Mediana ocorreu na região Sul ( 50,0 ), e a menor, na região Norte ( 40,2 ).
As demais medianas foram : 48,0 na região Nordeste, 47,4 na região Sudeste e 45,0 na região Centro-Oeste.
A nota Máxima do Brasil, como um todo, foi 96,5, sendo obtida por, pelo menos, um aluno na região Centro-Oeste.
As demais notas máximas foram : 84,6 na região Norte, 92,3 na região Nordeste, 92,9 na região Sudeste e 93,0 na região Sul.
A nota Mínima foi zero em todas as regiões, com exceção do 4,2 da região Norte.
Observa -se que existe diferença estatisticamente significativa da maior Média das notas do Componente de Conhecimento Específico da região Sul ( 49,9 ) em relação às demais regiões.
Assim como os Gráficos 6.1 e 6.2, o Gráfico 6.3, apresentado a seguir, permite uma avaliação do desempenho de concluintes de Arquitetura e Urbanismo em relação ao Componente de Conhecimento Específico com um histograma da distribuição das notas correspondentes.
Esta também é uma distribuição unimodal, e o grupo modal é o ( 40 ; 50 ], coincidindo com o grupo modal da prova como um todo e abaixo do grupo modal para a Formação Geral.
A Tabela 6.6 apresenta uma comparação dos resultados em relação à Categoria Administrativa, à Organização Acadêmica e à Modalidade de Ensino ( quando pertinente ), agora levando em conta o desempenho de alunos do Componente de Conhecimento Específico da prova.
No que se refere à Organização Acadêmica, a maior Média foi dos CEFET/IFET ( 51,1 ), seguidos pelas Universidades ( 49,5 ), vindo a seguir a dos Centros Universitários ( 44,3 ) e depois, das Faculdades ( 42,5 ).
O maior Desvio padrão, e acima do valor para o Brasil, como um todo, foi o das Universidades ( 16,2 ).
As Universidades também obtiveram a maior nota Máxima ( 96,5 ).
Os Centros Universitários obtiveram nota Máxima igual a 90,3, as Faculdades obtiveram 88,1, e finalmente os CEFET/IFET, 79,7.
As medianas foram : 50,3 dos CEFET/IFET, a maior delas, 50,1 nas Universidades, 44,1 nos Centros Universitários e 42,5 nas Faculdades, a menor.
A nota Mínima foi zero para todas as Organizações Acadêmicas, com exceção dos CEFET/IFET com 21,9.
Observa -se que existem diferenças estatisticamente significativas ao nível de 95 \% no Componente de Conhecimento Específico entre as Organizações Acadêmicas com menores médias, Faculdades e Centros Universitários.
Já a diferença entre a Média das Universidades e a dos CEFET/IFET não é estatisticamente significativa.
Quanto à Categoria Administrativa, observa -se um comportamento semelhante àquele da parte de Formação Geral e à prova como um todo, ou seja, existe diferença estatisticamente significativa entre as médias das IES Públicas ( 57,8 ) e IES Privadas ( 45,0 ).
Neste caso também, a maior Média foi obtida por alunos de IES Públicas de ensino. 6.2 ANÁLISE DAS QUESTÕES OBJETIVAS Esta seção apresenta estatísticas selecionadas e histogramas das Questões Objetivas de Formação Geral ( 6.2.1 ) e de Conhecimentos Específicos ( 6.2.2 ).
São também apresentadas e comparadas as médias das subpopulações caracterizadas por Grande Região. 6.2.1 Componente de Formação Geral A Tabela 6.7 apresenta as Estatísticas Básicas relativas às oito questões objetivas do componente da prova que abrange a Formação Geral dos estudantes.
A Média do Brasil foi 53,4.
A menor Média foi encontrada na região Norte ( 47,9 ), e a maior, na região Sul ( 56,1 ).
As demais médias foram : 54,5 na região Nordeste, 52,7 na região Sudeste e 52,6 na região Centro-Oeste.
O Desvio padrão do Brasil foi 20,5, sendo o maior Desvio padrão encontrado na região Nordeste ( 21,2 ), e o menor, na região Sul ( 19,8 ).
Os demais desvios foram : 20,1 na região Norte, 20,2 na região Sudeste e 21,1 na região Centro-Oeste.
As medianas do Brasil como um todo e de quatro das cinco regiões foi 50,0 - a única região com Mediana diferente foi a Sul, com 62,5.
A nota Máxima 100,0 foi alcançada em todas as Grandes Regiões.
As notas mínimas ( 0,0 ) foram iguais para todas as regiões.
A Tabela 6.8 apresenta o Índice de Facilidade e o Índice de Discriminação ( ponto-bisserial ) para cada uma das questões objetivas do Componente de Formação Geral.
Quanto ao Índice de Facilidade, foram usadas as seguintes cores para diferenciar o nível de dificuldade da questão :  Azul para as questões classificadas com índice Muito fácil ( \&gt ; =0,86 ), verde para as questões classificadas com índice Fácil ( 0,61 a 0,85 ), amarelo para as questões classificadas com Médio ( 0,41 a 0,60 ), vermelho para as questões classificadas com Difícil ( 0,16 a 0,40 ) e roxo para as questões classificadas com Muito difícil ( \&lt ; =0,15 ).
Já quanto ao Índice de Discriminação, foram usadas as seguintes cores para qualificar a questão : As questões objetivas do Componente de Formação Geral, segundo o Índice de Facilidade, foram assim avaliadas : das oito questões, nenhuma teve o Índice de Facilidade classificado como Muito difícil e uma questão foi tida como Muito Fácil, por ter índice de acertos 0,88.
Três questões foram consideradas com índice de dificuldade Médio, situando -se no intervalo entre 0,41 e 0,60 do Índice de Facilidade, ou seja, houve entre 41,0 \% e 60,0 \% de acertos, enquanto duas questões foram classificadas na categoria Difícil, situando -se no intervalo entre 0,16 e 0,40.
Por fim, duas questões apresentaram entre 61,0 \% e 85,0 \% de acertos, razão pela qual foram classificadas como Fácil.
O Índice de Facilidade variou de 0,19 a 0,88, e o de Discriminação, de 0,39 a 0,50.
A Tabela 6.9 apresenta a distribuição das questões considerando -se simultaneamente a classificação dos dois índices.
As seis questões com Índice de Discriminação Muito bom figuraram entre três níveis de dificuldade Médio, Fácil, e Muito Fácil : três classificadas na categoria Médio ( questões 2, 7 e 8 ) do Índice de Facilidade, duas na categoria Fácil ( as questões 4 e 5 ) e uma na categoria Muito Fácil ( questão 6 ).
Em particular, as questões 4 e 5 foram as que apresentaram o maior poder discriminatório, com índice 0,50, e foram consideradas Fácil em termos de facilidade, com uma proporção de, respectivamente, 0,71 e 0,64 de acertos.
O máximo de acertos foi alcançado pela questão 6 com um Índice de Facilidade de 0,88.
As questões 1 e 3 foram consideradas Difícil, com índices de facilidade de 0,19 e 0,36, respectivamente.
Já os índices de discriminação foram Bom para ambas as questões.  As questões classificadas com índice Fraco receberam a cor vermelho ( \&lt ; =0,19 ), as classificadas com Médio receberam a cor amarelo ( 0,20 a 0,29 ), as classificadas com Bom receberam a cor verde ( 0,30 a 0,39 ) e as classificadas com Muito bom ( \&gt ; =0,40 ) receberam a cor azul.
As questões objetivas do Componente de Formação Geral, segundo o Índice de Facilidade, foram assim avaliadas : das oito questões, nenhuma teve o Índice de Facilidade classificado como Muito difícil e uma questão foi tida como Muito Fácil, por ter índice de acertos 0,88.
Três questões foram consideradas com índice de dificuldade Médio, situando -se no intervalo entre 0,41 e 0,60 do Índice de Facilidade, ou seja, houve entre 41,0 \% e 60,0 \% de acertos, enquanto duas questões foram classificadas na categoria Difícil, situando -se no intervalo entre 0,16 e 0,40.
Por fim, duas questões apresentaram entre 61,0 \% e 85,0 \% de acertos, razão pela qual foram classificadas como Fácil.
O Índice de Facilidade variou de 0,19 a 0,88, e o de Discriminação, de 0,39 a 0,50.
A Tabela 6.9 apresenta a distribuição das questões considerando -se simultaneamente a classificação dos dois índices.
As seis questões com Índice de Discriminação Muito bom figuraram entre três níveis de dificuldade Médio, Fácil, e Muito Fácil : três classificadas na categoria Médio ( questões 2, 7 e 8 ) do Índice de Facilidade, duas na categoria Fácil ( as questões 4 e 5 ) e uma na categoria Muito Fácil ( questão 6 ).
Em particular, as questões 4 e 5 foram as que apresentaram o maior poder discriminatório, com índice 0,50, e foram consideradas Fácil em termos de facilidade, com uma proporção de, respectivamente, 0,71 e 0,64 de acertos.
O máximo de acertos foi alcançado pela questão 6 com um Índice de Facilidade de 0,88.
As questões 1 e 3 foram consideradas Difícil, com índices de facilidade de 0,19 e 0,36, respectivamente.
Já os índices de discriminação foram Bom para ambas as questões.
O Gráfico 6.4, para exemplificar, analisa o comportamento da questão de número 4 de Formação Geral.
Trata -se de uma questão considerada Fácil, com relação à facilidade e uma das que obtiveram o maior Índice de Discriminação dessa parte da prova.
Neste gráfico, cada uma das cinco curvas representa o percentual de respostas em determinada alternativa da questão, em função do número de acertos dos estudantes nessa parte da prova ( Formação Geral/Múltipla Escolha ), antes de possíveis eliminações pelo critério do ponto-bisserial.
A curva em verde corresponde à alternativa B, a alternativa correta para esta questão.
Observa -se que entre os estudantes com menor número de acertos, nessa parte do exame, a situação mais frequente foi a escolha da alternativa E ( em vermelho ), incorreta.
Por exemplo, entre os estudantes que acertaram três questões, mais de metade ( 53,8 \% ) escolheu a alternativa correta B ( em verde ), 19,0 \% escolheram a alternativa E, 13,7 \%, a D ( roxo ), 9,7 \%, a alternativa A ( azul ) e 3,8 \% escolheram a alternativa C ( preto ).
Entre os que acertaram três respostas nas questões de múltipla escolha de Formação Geral, 0,0 \% deixou a questão em branco e 0,1 \% marcou mais de uma alternativa, invalidando a questão.
À medida que o número de acertos aumenta, indicando desempenho melhor nessa parte da prova, aumenta concomitantemente a proporção de estudantes que selecionaram a alternativa correta B, atingindo 100 \% para os estudantes com 8 acertos.
Essa análise permite verificar como a questão discriminou os grupos de desempenho, justificando o alto índice obtido na questão.
Cumpre notar que não é possível inferir deste gráfico nem o índice de facilidade ( que seria uma média da proporção ponderada pela quantidade de alunos com cada uma das notas ), nem o índice de discriminação ponto-bisserial, por razão equivalente.
No caso extremo no qual a grande concentração dos acertos dos alunos fosse abaixo de 4, o índice de facilidade seria obrigatoriamente abaixo de 40 \% ( neste exemplo ).
Caso a concentração fosse em 6 acertos ou mais, o índice seria obrigatoriamente acima de 60 \%.
Os gráficos relativos às demais questões de Formação Geral constam do Anexo I. 6.2.2 Componente de Conhecimento Específico A Tabela 6.10 apresenta as Estatísticas Básicas em relação às questões objetivas do Componente de Conhecimento Específico da prova de Arquitetura e Urbanismo por Grande Região.
Nessa parte da prova nenhuma questão foi anulada pela Comissão Assessora da Área.
Além disso, como será discutido ainda nessa seção, quatro questões objetivas do Componente de Conhecimentos Específicos não foram usadas no cômputo das notas por terem sido descartadas pelo critério do ponto-bisserial.
Assim, as notas foram calculadas com base em 23 das 27 questões objetivas de Conhecimentos Específicos.
A Média do Brasil deste componente foi de 49,2.
A maior Média foi observada na região Sul ( 52,0 ), e a menor, na região Norte ( 42,7 ).
O Desvio padrão de todo o Brasil foi 16,8, sendo o maior Desvio padrão encontrado nas regiões Nordeste e Centro-Oeste ( 17,4 ), e o menor, na região Norte ( 15,6 ).
A Mediana de todo o Brasil foi 47,8.
Nas regiões Nordeste, Sudeste e Centro-Oeste, a Mediana foi a mesma ( 47,8 ), na região Sul foi maior ( 52,2 ) e na região Norte foi menor ( 43,5 ).
A nota Máxima da prova foi 100,0, obtida nas questões objetivas do Componente de Conhecimento Específico, por, pelo menos, um aluno de três das regiões : Nordeste, Sul e Centro-Oeste.
A nota Máxima foi menor para a região Norte ( 91,3 ) e para a região Sudeste ( 95,7 ).
A nota Mínima foi zero em todas as regiões, com exceção da região Norte ( 4,3 ).
A Tabela 6.11 apresenta os Índices de Facilidade e Discriminação ( ponto-bisserial ) das questões objetivas do Componente de Conhecimento Específico para os estudantes de Arquitetura e Urbanismo.
Para facilitar a diferenciação das questões, usaremos as mesmas cores da Tabela 6.8 para as diferentes classificações dos Índices de Facilidade e de Discriminação.
Observa -se que nenhuma questão foi anulada pela CAA, e as análises dizem respeito às 27 questões válidas.
A partir do Índice de Facilidade obtido, pode -se concluir que pouco mais do que um terço das questões objetivas da prova foram consideradas Difícil : das 27 questões válidas, 11 foram classificadas como Difícil e nenhuma, como Muito difícil.
Cinco questões foram classificadas como Fácil, e outras 11 consideradas como Médio.
Já quanto aos índices de discriminação das questões objetivas do Componente de Conhecimento Específico da prova, tem -se como resultado a seguinte classificação : 12 das 27 questões foram consideradas boas, enquanto seis delas tiveram Índice de Discriminação Muito bom.
Assim, para dois terços das questões – 18 em 27 – os Índices de Discriminação foram Bom ou Muito bom.
Dentre as demais, cinco delas foram classificadas como Médio e outras quatro como Fraco, sendo nove, por conseguinte, a quantidade de questões nos dois patamares mais baixos de discriminação.
Constata -se, assim, que a prova – no que se refere ao Componente de Conhecimento Específico – possuía boa capacidade de discriminar entre aqueles que dominam ou não o conteúdo.
A Tabela 6.12 apresenta a distribuição das questões considerando -se simultaneamente a classificação dos dois índices.
Dentre as questões que alcançaram os maiores Índices de Discriminação, classificadas com Muito bom, uma delas, a de número 23, foi a de índice mais elevado ( 0,47 ), classificada como Fácil quanto ao Índice de Facilidade, com 66 \% dos estudantes marcando a opção correta.
A questão de número 20 foi a mais difícil dentre as 27 questões específicas válidas, com baixo Índice de Facilidade, apenas 16 \% de acertos.
Essa questão apresentou poder discriminatório muito baixo, 0,04, o que comprova ter sido a mais difícil para os estudantes.
Destaca -se, também, a questão 22, com Índice de Facilidade 0,18, o que, em termos percentuais, corresponde a 18 \% de estudantes que responderam acertadamente.
Já 0,18 foi o seu Índice de Discriminação.
Tais questões foram, portanto, pelo critério ponto-bisserial, consideradas inadequadas.
Por isso, as questões 20 e 22, além das questões 11 e 32 foram eliminadas do cômputo da nota final.
A título de exemplo das análises do comportamento das questões objetivas, o Gráfico 6.5 analisa a questão 12 do Componente de Conhecimento Específico.
Esta questão foi considerada pelas respostas dos estudantes avaliados na prova como uma questão com nível de facilidade Médio, 0,51, ou seja, 51 \% dos estudantes assinalaram acertadamente a opção E, correspondente ao gabarito.
Como já comentado, seu Índice de Discriminação foi igual a 0,44, classificado como Muito bom, e o segundo maior valor de discriminação.
Neste gráfico, cada uma das cinco curvas representa o percentual de respostas em determinada alternativa da questão 12, em função do número de acertos dos estudantes nessa parte da prova, antes de possíveis eliminações de questões pelo critério do ponto-bisserial.
A alternativa correta E, representada no gráfico pela curva em vermelho, foi escolhida em maiores proporções pelos alunos com desempenho melhor nessa parte da prova.
Já as alternativas incorretas, também denominadas distratores, foram selecionadas, principalmente, por aqueles com notas mais baixas.
Observa -se que a soma não é 100 \% por causa das questões não respondidas ou com mais de uma opção marcada.
Aqueles com nota zero, na sua quase totalidade, deixaram esta questão em branco ou marcaram mais de uma alternativa, comportamento considerado inválido.
A proporção de alunos que selecionou a resposta correta E aumenta gradativamente, chegando a atingir 100 \% para 27 acertos, enquanto a proporção dos que escolheram alternativas incorretas decai, a partir de quatro ou cinco acertos, como função do número de acertos nessa parte da prova.
Os gráficos relativos às demais questões do Conhecimento Específico constam do Anexo I. 6.3 ANÁLISE DAS QUESTÕES DISCURSIVAS Esta seção apresenta estatísticas selecionadas e histogramas das Questões Discursivas de Formação Geral ( 6.3.1 ) e Conhecimento Específico ( 6.3.2 ).
São também apresentadas e comparadas às médias de subpopulações caracterizadas por Grande Região. 6.3.1 Componente de Formação Geral As análises dos resultados de desempenho dos estudantes de Arquitetura e Urbanismo nas duas questões discursivas relativas à Formação Geral encontram -se na Tabela 6.13 e no Gráfico 6.6.
Na Tabela 6.13, observa -se que a nota Média nesse conjunto de questões foi acima da obtida nas objetivas.
Os estudantes, de todo o Brasil, obtiveram, em Formação Geral, Média 53,4 nas questões objetivas e 52,2 nas questões discursivas.
Pode -se notar também que o Desvio padrão nesse conjunto de questões foi um pouco maior do obtido nas objetivas, 20,5 nas questões objetivas e 24,6 nas questões discursivas.
A maior Média foi obtida na região Sul ( 56,1 ), e a menor, na região Norte ( 52,2 ).
O Gráfico 6.6 representa a distribuição das notas nas questões discursivas do Componente de Formação Geral.
A moda desta distribuição ocorre no intervalo ( 60 ; 70 ], com frequência um pouco maior do que a do intervalo ( 70 ; 80 ].
Destaca -se também o intervalo [ 0 ; 10 ] com distribuição próxima a 10 \% do total de notas, sendo que no intervalo [ 0 ; 10 ] inclui -se, além da nota zero, a frequência de alunos que deixaram as respostas a este tipo de questão em branco.
A Mediana de todo o Brasil, neste componente, foi 61,5.
A Mediana no Norte ( 57,0 ) e a do Centro-Oeste ( 60,0 ) foram menores que a nacional.
Nas demais regiões, foi maior : 62,0 no Nordeste e no Sudeste e 62,5 no Sul.
A nota Máxima foi 99,0, obtida por, pelo menos um estudante da região Nordeste.
Nas outras quatro das cinco Grandes Regiões, a nota Máxima foi 98,0.
A nota Mínima ( 0,0 ) foi a mesma em todas as regiões do Brasil.
O Gráfico 6.6 representa a distribuição das notas nas questões discursivas do Componente de Formação Geral.
A moda desta distribuição ocorre no intervalo ( 60 ; 70 ], com frequência um pouco maior do que a do intervalo ( 70 ; 80 ].
Destaca -se também o intervalo [ 0 ; 10 ] com distribuição próxima a 10 \% do total de notas, sendo que no intervalo [ 0 ; 10 ] inclui -se, além da nota zero, a frequência de alunos que deixaram as respostas a este tipo de questão em branco.
Na sequência, os resultados verificados para cada uma das questões discursivas de Formação Geral serão apresentados, estabelecendo -se relações com os conteúdos abordados em cada uma delas.
Os comentários da Banca de docentes corretores, a respeito do observado na correção das respostas dos estudantes, suas impressões e conclusões serão apresentados junto à análise de cada questão.
Cumpre esclarecer que, tendo em vista que as questões discursivas de Formação Geral são padronizadas, ou seja, constam de todas as provas, os comentários da Banca são os mesmos para todas as carreiras acadêmicas, sendo direcionados a todos os estudantes que participaram do Enade/2017.
A seguir, serão analisados os desempenhos da Área de Arquitetura e Urbanismo nas duas questões discursivas de Formação Geral do Enade/2017, comparando -se os resultados obtidos com comentários para cada questão. 6.3.1.1 Análise de Conteúdo da Questão Discursiva 1 do Componente de Formação Geral Os dados de Arquitetura e Urbanismo, obtidos a partir das respostas à questão 1, encontram -se na Tabela 6.14 e no Gráfico 6.7.
Nessa questão – de desempenho dentre as duas de Formação Geral – os alunos, de todo o Brasil, tiveram Média 55,0.
A maior Média para a questão 1 foi obtida na região Sudeste ( 55,5 ), e a menor, na região Norte ( 52,9 ).
Quanto à variabilidade das notas, o Desvio padrão de todo o Brasil foi 28,7.
O menor Desvio padrão foi obtido na região Sudeste ( 27,9 ), e o maior Desvio padrão foi obtido na região Nordeste ( 30,5 ).
As medianas de todas as regiões e do Brasil como um todo foram iguais ( 60,0 ), com exceção da região Norte com Mediana 55,0.
As notas máximas e mínimas da questão discursiva 1 foram, sem exceção, as mesmas para todas as regiões do Brasil, respectivamente, 100,0 e 0,0.
O Gráfico 6.7 mostra a distribuição das notas na questão discursiva 1 do Componente de Formação Geral.
Observa -se que a maior frequência ocorre no intervalo ( 70, 80 ], destacando -se, também, a frequência dos alunos que deixaram a resposta a essa questão em branco, em torno de 12 \% dos participantes. 6.3.1.2 Comentários sobre a correção de Conteúdo das respostas à Questão Discursiva 1 O enunciado era claro e indicava um tema de extrema relevância, qual seja, a incidência de doença grave integrante do rol das DST ( Doenças Sexualmente Transmissíveis ), que encaminha para a análise do direito fundamental à saúde, que tem \textbf{assento} constitucional ( CRFB, art. 6º. caput : São direitos sociais a educação, a saúde, a alimentação, o trabalho, a moradia, o transporte, o lazer, a segurança, a previdência social, a proteção à maternidade e à infância, a assistência aos desamparados, na forma desta Constituição ).
A questão focalizava, por meio de três textos, um tema social extremamente relevante, qual seja, a relação entre homem e mulher nas suas interações íntimas, no caso representadas pela vida sexual, e o reflexo das relações de poder veiculadas pela sociedade, que apontam, no mais das vezes, o predomínio do homem.
Tal posição de superioridade é justificada pelas convenções arraigadas que acarretam a educação dos homens e mulheres por padrões patriarcais, que redundam em posições machistas.
Um dos textos demandava uma reflexão sobre a realidade dos \textbf{casais} que convivem em relações estáveis ou fugazes e não têm o controle dos efeitos das relações sexuais, que abrangem a gravidez não planejada e, em decorrência do não cuidado, a inserção de doenças transmitidas sexualmente.
Esse panorama é refletido na pouca influência das mulheres sobre os homens quanto à proteção, normalmente representada pelo uso de preservativos nas relações intimas, e após a constatação da doença, no encaminhamento das pessoas para tratamento.
Como efeito colateral, está a transmissão de doenças para os fetos.
Outro texto trazia à tona o maior cuidado que as mulheres têm, em geral, com sua saúde, notadamente no momento em que se descobre grávida e que redunda na realização de exames obrigatórios de pré-natal, que podem facilmente detectar doenças e, caso encontradas, ter o seu tratamento encaminhado.
A par disso, destaca -se que a crise de recursos financeiros agravou a situação, que já recebe influxos sociais negativos, como a dificuldade do homem em utilizar preservativos, realizar exames e procurar tratamentos quando adoece.
Como muito divulgado pelas mídias, a crise financeira do país repercutiu diretamente na prestação de serviços de saúde em todas as esferas : distrital, federal, estadual e municipal.
Assim, agravou -se o problema da falta de medicamentos essenciais para o tratamento da sífilis e, certamente, de outras doenças.
Como se verifica, a temática indicada no enunciado, representada pelos textos, pertence à realidade conhecida amplamente pelos estudantes, pois também é divulgada por todos os meios de comunicação.
Trata -se de tema geral que, no entanto, normalmente, não se inclui nas grades curriculares de cursos superiores que não sejam da área de Saúde.
Apesar disso, a temática é de conhecimento geral e relacionada aos direitos humanos, esses sim, ponto de integração com as diretrizes curriculares.
A linguagem utilizada era clara, indicando a existência dos fatos vinculados a questões relevantes e à saúde pública.
No entanto, como os textos trazem uma ampla gama de possibilidades de interpretação, ocorreram algumas compreensões diversas das esperadas.
Além disso, ressalta -se que a amplitude dos textos disponibilizados no enunciado fornece elementos que poderiam ser simplesmente reproduzidos pelos concluintes nas suas respostas.
O enunciado da questão encaminha a resposta para dois aspectos.
O estudante deveria abordar conteúdo pertinente à vulnerabilidade das mulheres às DST e o papel social do homem em relação à prevenção dessas doenças.
Para responder a esses aspectos, um número expressivo de concluintes utilizou somente os elementos dos textos motivadores, sem nada acrescentar.
A resposta deveria conter também uma segunda parte, a indicação de duas ações voltadas para o público masculino, a serem adotadas no âmbito das políticas públicas de saúde ou de educação, para a redução do problema.
Verificou -se que o desempenho dos estudantes nessa questão foi de médio para bom.
Ao se analisarem as respostas, constata -se que muitos foram os que apresentaram os tópicos presentes no padrão de resposta, amplo o suficiente para albergar as respostas apresentadas pelos concluintes.
Tal resultado pode ser justificado, tendo em vista tratar -se de tema bastante difundido nas diversas mídias, por ter sido apresentado de forma ampla pelos textos motivadores presentes no enunciado.
Por outro lado, o debate no ambiente universitário sobre as relações de poder entre homens e mulheres tende a ser marcante pelo movimento de empoderamento feminino, liderado por organizações sociais com ampla atuação nos estabelecimentos de nível superior.
De forma recorrente, as relações patriarcais entre homens e mulheres foram vinculadas ao modelo machista vigente.
Essas condições, segundo os concluintes, explicam a incidência da doença e a situação de vulnerabilidade das mulheres, bem como a função social dos homens.
Outras respostas apontaram o poder público, nas áreas de saúde e/ou educação, como órgãos responsáveis pela solução do problema.
Para os que se fixaram na área da saúde, foram frequentes a indicação da necessidade de mais medicamentos, vacinas, educação sexual, médicos e postos de atendimentos.
Alguns, no entanto, apenas indicaram ações individuais para o homem, e outros, ações individuais para as mulheres.
Também figuraram respostas que apresentavam ações da igreja como solução, além de outras que faziam referências a campanhas no seio das famílias. 6.3.1.3 Análise de Conteúdo da Questão Discursiva 2 do Componente de Formação Geral A Tabela 6.15 mostra que o desempenho médio dos estudantes na questão discursiva 2 ( média 55,1 ) foi bem próximo ao obtido na questão discursiva 1 ( média 55,0 ).
A região Norte foi aquela cuja Média, também nessa questão, foi menor ( 51,3 ), e a de maior Média foi a região Sul ( 56,6 ).
Quanto à variabilidade das notas, o Desvio padrão de todo o Brasil foi 32,8, superior ao obtido na questão discursiva 1 ( 28,7 ).
O maior desvio nessa questão foi obtido na região Nordeste ( 34,4 ), enquanto o menor foi obtido nas regiões Sudeste e Sul ( 32,0 ).
A Mediana de todo o Brasil foi 70,0, a mesma de todas as regiões, a menos da região Norte, cuja Mediana foi 50,0.
As notas máximas ( 100,0 ) e as notas mínimas ( 0,0 ) também foram as mesmas em todas as regiões, sem exceção.
O Gráfico 6.8 mostra a distribuição das notas na questão discursiva 2 do Componente de Formação Geral.
Observa -se que a maior frequência corresponde aos alunos que obtiveram nota dentro do intervalo ( 70 ; 80 ].
Destacam -se, também, os alunos que deixaram a resposta à questão em branco e os alunos que obtiveram nota dentro do intervalo ( 90 ; 100 ], caracterizando -se como máximos locais.
Nota -se, ainda, que as notas ficaram um pouco menos concentradas em comparação à questão discursiva de número 1.
Isto pode ser constatado também pela comparação do desvio padrão das notas da questão discursiva 2 ( 32,8 ) e da questão discursiva 1 ( 28,7 ). 6.3.1.4 Comentários sobre a correção de Conteúdo das respostas à Questão Discursiva 2 O enunciado era claro e abordava um tema de extrema relevância : a questão de gênero.
Trata -se de uma proposta polêmica, pois abrange a análise da inserção social de minoria vinculada à questão de gênero.
O enunciado solicitava que o concluinte discorresse sobre a importância do nome para as pessoas transgêneras e propusesse uma medida, no âmbito das políticas públicas, para facilitar o acesso dessas pessoas à cidadania.
A respeito da primeira parte, a importância do nome, esperava -se que o estudante abordasse as dificuldades que o uso do nome civil ( nome de registro, nome de batismo, nome da certidão de nascimento, por exemplo ) pode acarretar para as pessoas transgêneras, além de demonstrar conhecimento sobre a importância do nome social, direito já reconhecido através, principalmente, de decisões judiciais e atos administrativos, como decretos e portarias que dispõem sobre o uso do nome social.
Em atendimento à segunda parte do comando, o participante deveria propor uma política pública associada ao direito à cidadania das pessoas transgêneras.
Sobre esse aspecto, esperava -se a compreensão do significado de política pública, ou seja, propostas de âmbito governamental, e não, iniciativas de Organizações Sociais, igrejas ou grupos sociais de qualquer tipo.
O nascimento das pessoas, no Brasil, é objeto de inscrição em \textbf{cartórios} de Registro Civil das Pessoas Naturais, órgão responsável por estabelecer a certificação pública dos dados de todos os indivíduos que não somente nascem, mas casam, eventualmente se divorciam e vêm a óbito.
Daí o nome civil.
Para flexibilizar a correção, outras nomenclaturas foram consideradas, a saber : identificação de nascimento, nome de batismo, nome de registro, nome próprio, nome verdadeiro, que diferenciam o nome civil daquele que o indivíduo deseja utilizar para traduzir a sua opção pessoal de gênero, que se convencionou denominar nome social.
Quanto ao entendimento do termo transgênero, deve -se destacar que o movimento de mudança do gênero é contemporâneo e não está necessariamente vinculado à mudança de sexo.
Existem várias hipóteses em jogo.
Uma é a que envolve cirurgias de mudança de sexo ( a menina extirpa os seus seios e inclui órgão genital por intervenção cirúrgica ; o menino inclui seios e extirpa o seu órgão genital e inclui o órgão feminino também por intervenção cirúrgica ).
Há outro movimento que envolve apenas as vestimentas e o linguajar corporal, além dos relacionamentos ( as meninas assumem o trajar, o andar, os gostos dos meninos e mantêm relacionamentos afetivos com meninas ; os meninos assumem o trajar, o andar, os gostos das meninas e mantêm relacionamentos afetivos com meninos ).
Em relação à discussão da primeira solicitação do comando da questão, destaca -se que todos têm nome civil, que foi escolhido pelos seus pais ou responsáveis e, na falta deles, por autorização do Juiz de Direito, pautado no sexo detectado no nascimento.
Já o nome social é aquele que indica a preferência do indivíduo na sua identificação nas relações sociais e que destoa do seu gênero de registro civil.
Ressalte -se que o Brasil não possui ainda legislação específica sobre pessoas transgêneras.
Observa -se um relativo progresso no reconhecimento de direitos, através, principalmente, de decisões judiciais e atos administrativos, como decretos e portarias que dispõem sobre o uso do nome social.
Sem uma lei que defina os procedimentos da alteração dos documentos para pessoas transexuais, essa parcela da população LGBT ( lésbicas, gays, bissexuais, transexuais e transgêneros ) é obrigada a procurar na Justiça o reconhecimento de sua identidade, em processos que podem ser longos e que dependem do olhar de determinados, profissionais, como médicos, psicólogos, assistentes sociais, bem como do entendimento dos \textbf{advogados}, defensores públicos, juízes e do Judiciário, como um todo.
Atualmente, tramita na Comissão de Direitos Humanos da Câmara dos Deputados o Projeto de Lei João Nery ( 5002/2013 ), dos deputados Jean Willys ( PSOL-RJ ) e Erika Kokay ( PT-DF ), que determina que o reconhecimento da identidade de gênero seja um direito do cidadão.
O projeto recebeu o nome do primeiro trans-homem operado no Brasil.
Grande parte dos respondentes considerou a incongruência entre a imagem corporal, representada pelo gênero assumido pelos trans, e o nome que apresenta nos seus documentos como causa de sérios constrangimentos e embaraços e as dificuldades para que esses sujeitos acessem direitos básicos, como saúde, educação e empregabilidade em uma sociedade heteronormativa e preconceituosa.
Muitos demonstraram reconhecer que a identidade e expressão de gênero, através do nome social, pressupõe que sejam respeitadas a dignidade da pessoa humana, a intimidade, a liberdade e o direito à cidadania plena.
No entanto, também foram encontradas, em quantidades significativas, respostas que demonstravam a existência de uma interpretação social conservadora que não admite a mudança de sexo nem do nome social.
Existe um acentuado preconceito contra todos os movimentos que buscam identificação social diversa da civil.
Isso inclui os variados grupos : gays, travestis, transgêneros, dentre outros.
Já em relação ao segundo tópico a ser abordado pelo concluinte, é notório que as pessoas trans encontram -se destituídas de direitos no campo das políticas públicas, como saúde, habitação, emprego, educação, justiça, dentre outros.
Nesse sentido, uma grande diversidade de medidas, no âmbito das políticas públicas, poderia ser proposta.
Ações voltadas para a capacitação de profissionais da educação, saúde, segurança pública, dentre outras áreas, para melhor atender às demandas e especificidades desse público alvo seriam necessárias.
Ainda em relação à educação, seriam também necessárias ações que contribuíssem para diminuir a evasão, pois o ambiente escolar é difícil para uma pessoa transgênera, por conta das situações discriminatórias e preconceituosas vindas dos colegas de classe, e mesmo das/dos professores e/ou de outros profissionais da área, não podem ser mais um \textbf{complicador} nesse processo.
Ainda no que tange ao assunto, percebe -se que, ao serem excluídas dos bancos escolares, as oportunidades no mercado de trabalho formal praticamente inexistem.
Um outro aspecto que se evidencia no cotidiano desse segmento da sociedade diz respeito às situações de violência sofridas pelas pessoas transgêneras.
Nesse sentido, algumas propostas foram elencadas para minimizar o problema e combater a violência.
Políticas envolvendo melhor formação de funcionários das polícias civil e militar, por exemplo, somariam esforços na qualificação das abordagens policiais e dos atendimentos prestados pelas forças de segurança, pois o respeito à orientação sexual e à identidade e expressão de gênero é, antes de tudo, o respeito à dignidade, à cidadania e à própria democracia.
Nesse sentido, revela -se essencial que, ao sofrerem agressões, sejam elas físicas, psicológicas ou morais no espaço público ou domiciliar, as pessoas trans possam ser atendidas, com respeito inclusive ao nome social, nas delegacias comuns e, inclusive, nas Delegacias Especiais de Atendimento à Mulher ( DEAM ), pois esse é o órgão institucional que possivelmente detém as qualificações necessárias para o atendimento ao gênero feminino ao qual as travestis e mulheres trans efetivamente pertencem.
Também ocorre uma latente inacessibilidade das pessoas trans à assistência jurídica.
Nesse sentido, é necessário implementar projetos de capacitação para profissionais do poder judiciário, como um todo, de modo a prestarem a assistência jurídica gratuita na defesa intransigente dos direitos sociais, civis, trabalhistas e políticos dessa população.
Dessa maneira, a Justiça agiria de forma mais acessível e sensível para o segmento, que terá suas especificidades observadas a começar pelo devido respeito do direito ao uso do nome social.
Nesse âmbito, o mais comum foi propor a adoção de medidas punitivas para quem viola o direito à autodeterminação de gênero e a elaboração de leis que garantam a mudança do nome e assegurem outros direitos para as pessoas trans.
Propostas envolvendo a formação de profissionais, como as mencionadas aqui, e de conscientização da população foram as mais frequentes.
Muitas outras enfatizavam o direito ao uso respeitado do nome social e a garantia de alteração do registro civil, independentemente da cirurgia de redesignação sexual ou transgenitalização.
Muitos questionaram o fato de não bastar a autodeclaração no processo de alteração de nome, enfatizando que o nome social é uma questão de princípio de dignidade e cidadania.
Apesar de pouco mencionada nas respostas corrigidas, políticas envolvendo a promoção da saúde são muito necessárias.
Cabe aos gestores de saúde ações direcionadas, não somente à prevenção e tratamento do HIV/Aids, mas também políticas de saúde que contemplem as especificidades dessa população, tais como o Ministério da Saúde, especialmente o SUS ( Sistema Único de Saúde ), incluir, de fato, o atendimento para pessoas trans para acompanhamento da terapia hormonal, que muitas demandam, respeitando -se o direito de utilização do nome social.
De modo geral, são muitas as políticas públicas a que precisam ser implantadas.
Urge a necessidade de uma agenda mais compromissada por parte dos gestores do nosso país com essa comunidade, que ainda é muito discriminada.
Logo, tais políticas devem contemplar suas necessidades básicas como o direito de acesso aos estudos, à capacitação para o mercado de trabalho, à profissionalização, ao acesso a bens e serviços de qualidade em saúde, habitação, segurança, cultura e assistência social e, certamente, o respeito ao uso do nome social, o que também contribui para essa inclusão. mercado de trabalho, à profissionalização, ao acesso a bens e serviços de qualidade em saúde, habitação, segurança, cultura e assistência social e, certamente, o respeito ao uso do nome social, o que também contribui para essa inclusão.
Assim, as opções apresentadas no padrão de resposta funcionaram como exemplos de políticas públicas possíveis, já que seria difícil apresentar uma lista suficientemente ampla para abranger a grande variedade de respostas encontradas.
O desempenho dos estudantes nessa questão foi médio.
As respostas, de modo geral, indicaram a necessidade de uma atualização por parte dos concluintes quanto a temas essenciais e, no caso da presente questão, a análise sobre aspectos fundamentais da questão de gênero.
Muitos, possivelmente motivados pelos textos apresentados no enunciado, restringiram -se a discorrer sobre preconceito de gênero ou mesmo sobre o preconceito de modo mais geral, sem responder ao que foi solicitado no comando da questão.
Outros apresentaram discursos religiosos ou ideológicos, também sem tratarem do que foi solicitado. 6.3.1.5 Análise de Língua Portuguesa das Questões Discursivas do Componente de Formação Geral Os dados de Arquitetura e Urbanismo, obtidos a partir das respostas às questões discursivas do Componente de Formação Geral, no que tange à Língua Portuguesa, encontram -se na Tabela 6.16 e no Gráfico 6.9.
Nesse aspecto, os alunos, de todo o Brasil, obtiveram Média 56,3.
A maior Média com respeito à Língua Portuguesa foi obtida na região Sul ( 57,6 ), e a menor, na região Norte ( 53,0 ).
Quanto à variabilidade das notas, o Desvio padrão de todo o Brasil foi 22,5.
O menor Desvio padrão foi obtido na região Sudeste ( 21,5 ) e o maior Desvio padrão foi obtido na região Nordeste ( 24,3 ).
A Mediana das notas de Língua Portuguesa foi 62,5 para o Brasil, como um todo, a mesma obtida em três regiões : Nordeste, Sudeste e Centro-Oeste.
A Mediana da região Norte foi 57,5 e a da região Sul, 65,0.
A nota Máxima para todo o Brasil foi 97,5, com, pelo menos, um aluno tirando essa nota nas regiões Nordeste e Centro-Oeste.
Na região Sul, a nota Máxima foi 95,0, enquanto nas regiões Sudeste e Norte foi 92,5.
Já a nota Mínima foi zero em todas as regiões do país.
O Gráfico 6.9 mostra a distribuição das notas de Língua Portuguesa do Componente de Formação Geral.
Observa -se que a maior frequência ( quase 30 \% ) corresponde aos alunos que obtiveram nota no intervalo ( 60 ; 70 ].
Destacam -se, também, os alunos que deixaram ambas as respostas às questões discursivas de Formação Geral em branco, aproximadamente 8 \% do total, caracterizando -se como máximo local. 6.3.1.6 Comentários sobre a correção das respostas de Formação Geral com respeito à Língua Portuguesa Os enunciados apresentados em 2017 não explicitaram a exigência de elaboração de um texto “ dissertativo ”, como nos anos anteriores.
O participante deveria fazer uma exposição de seus conhecimentos sobre cada um dos assuntos e estruturar seus textos de acordo com as características do registro formal adequado à situação comunicativa – avaliação de conhecimentos.
Essa configuração determina exigências quanto aos seguintes aspectos : adequação da seleção vocabular, desenvolvimento do conteúdo, organização lógica das ideias, estruturação sintática dos períodos, utilização de procedimentos de encadeamento textual e de referenciação, obediência às exigências morfossintáticas próprias da modalidade escrita da norma-padrão, respeito às regras ortográficas e às regras de acentuação gráfica.
O padrão de resposta utilizado na avaliação das questões 1 e 2 considerou os aspectos relevantes ao bom desempenho linguístico como competências distintas, de modo a permitir um mapeamento detalhado do domínio dos recursos disponíveis na Língua Portuguesa para a comunicação escrita formal.
Com base nesse objetivo, foram avaliados os seguintes aspectos : ( a ) estruturação textual condizente com o gênero solicitado e o modo de organização textual expositivo adequado ao gênero – essa competência envolve : • estruturação sintática condizente com o padrão da modalidade escrita formal da língua portuguesa de modo a garantir a clareza necessária ; • distribuição do conteúdo do texto em parágrafos, de modo a garantir a sua organização temática ; • utilização de operadores discursivos que contribuam para a progressão temática do texto, estabelecendo relações lógicas entre as ideias apresentadas, tanto do ponto de vista intrafrasal, como do interfrasal ; • utilização de procedimentos de referenciação lexical e pronominal que permitam a retomada de referentes textuais ; • utilização de sinais de pontuação que contribuam para a organização lógica da frase e do texto ; • inteligibilidade relacionada ao atendimento das exigências de estruturação textual.
Espera -se, portanto, que o participante recorra a procedimentos linguístico-discursivos para organizar seu texto, permitindo o encadeamento lógico entre suas partes de forma a garantir a progressão e a coerência textuais.
Isso significa que os seguintes procedimentos foram considerados inadequados, de acordo com o padrão de resposta proposto :  elaboração de frases fragmentadas que comprometam a estrutura lógico-gramatical do texto ;  sequência justaposta de ideias sem encaixamentos sintáticos, reproduzindo hábitos da oralidade ;  elaboração de frase com apenas oração subordinada, sem oração principal ;  emprego equivocado do conector ( preposição, conjunção, pronome relativo, alguns advérbios e locuções adverbiais ) comprometendo a expressão da relação lógica entre duas ideias, com prejuízo da clareza do texto ;  emprego do pronome relativo sem a preposição, quando obrigatória ;  repetição ou substituição inadequada de palavras sem o emprego dos recursos oferecidos pela língua ( pronome, advérbio, artigo, sinônimo ) ;  emprego inadequado dos pronomes relativos “ cujo(a) ” e “ onde ” ;  utilização inadequada dos sinais de pontuação, comprometendo a clareza textual ;  inteligibilidade relacionada ao atendimento das exigências de estruturação textual. ( b ) respeito às convenções ortográficas da norma-padrão da Língua Portuguesa – essa competência envolve o domínio das regras de acentuação gráfica e da grafia padrão das palavras ( com ausência de abreviaturas próprias da linguagem da internet ), de acordo com as convenções estabelecidas pela legislação em vigor e consubstanciadas no Vocabulário Ortográfico da Língua Portuguesa, editado pela Academia Brasileira de Letras. ( Vale notar que, nessa edição de 2017, não foi aceita a legislação anterior, no caso das regras relativas ao uso do hífen e da acentuação gráfica ).
Espera -se que o participante :  grafe corretamente as palavras ;  respeite as regras de acentuação gráfica ;  empregue maiúsculas em início de frase, em nomes próprios de pessoas, lugares ou instituições ;  grafe as siglas com letras maiúsculas ou apenas com a primeira letra maiúscula quando formarem uma palavra ;  evite abreviações como p/, vc, tb, pq, tá, né, usadas muitas vezes em escrita informal e na internet ;  obedeça às regras de separação de sílabas no final da linha. ( c ) domínio dos diferentes aspectos morfossintáticos próprios da modalidade escrita formal da norma-padrão da Língua Portuguesa – essa competência envolve : a concordância nominal, a concordância verbal, a regência nominal, a regência verbal, a flexão nominal, a flexão verbal, a correlação entre os tempos verbais, a colocação pronominal e a utilização de sinais de pontuação que contribuam para a organização lógica da frase e do texto.
Espera -se que o participante :  flexione o verbo para estabelecer concordância de número com o sujeito da frase ;  flexione o artigo, o adjetivo e o pronome para concordar em número e em gênero com o substantivo a que se referem ;  observe a regência nominal e a verbal, utilizando a preposição adequada depois de um substantivo, um verbo ou um adjetivo ;  empregue adequadamente o acento grave indicador da crase entre uma preposição e um artigo ( a+a ) ;  obedeça às regras de colocação pronominal ( próclise e ênclise ), distintas dos hábitos da oralidade ou da escrita informal ( exigência de próclise com termo atrator, não exigência de que o pronome oblíquo se ligue ao verbo auxiliar por meio de hífen ) ;  flexione adequadamente verbos, substantivos, adjetivos e pronomes no que diz respeito à expressão das categorias gramaticais ;  flexione os verbos para expressar a correlação de modo e tempo nas estruturas subordinadas.
Com base nesses critérios, foram considerados como desvios de caráter morfossintático, e não como desvios ortográficos, as alterações que envolvem mudança de classe gramatical ou de forma flexional do verbo :  eliminação da marca de infinitivo ( -r- ) e substituição por acento agudo ou ausência total de marca do infinitivo ;  confusão entre “ ão ” e “ am ” nas formas verbais ;  confusão entre “ há ” e “ a ” ;  uso de hífen para separar pronome átono – tanto uso indevido quanto omissão ( exemplo : “ esperasse ”, em lugar de “ espera -se ” ; “ falar -mos ”, no lugar de “ falarmos ” ;  “ esta ” ( no lugar de “ está ” ) ; “ mais ” ( no lugar de “ mas ” ) ; “ e ” ( no lugar de “ é ” ).  verbos “ ter ” e “ vir ” que, na terceira pessoa do plural, não apresentarem o acento circunflexo, serão considerados como desvio de concordância, em aspectos morfossintáticos. ( d ) seleção vocabular adequada à modalidade escrita formal da Língua Portuguesa, exigida pela situação comunicativa – essa competência envolve a precisão na seleção/utilização do vocabulário relacionado à temática solicitada pela questão ; a ausência de marcas da oralidade, como termos de sentido muito genérico ( “ coisa ”, “ negócio ”, “ você ” ) e termos de registros mais informais ( como gírias, jargões, frases feitas, ditados populares, termos regionais ).
Assim, espera -se que o participante respeite a adequação vocabular não usando gírias ou expressões coloquiais, evite repetição desnecessária de palavras e utilize um vocabulário mais formal, como solicitado por um texto dissertativo. \textbf{Observações} : - A inteligibilidade relacionada ao atendimento das exigências de estruturação textual foi avaliada na segunda competência ( relativa aos aspectos textuais ). - Os problemas de coerência textual provocados por uso indevido do vocabulário foram avaliados na terceira competência ( relativa aos aspectos morfossintáticos e vocabulares ). - Não foram considerados como desvios morfossintáticos os problemas de caligrafia ( -a/ -o, -s, -r ). - Não foram considerados os textos grafados integralmente em caixa alta, já que não seria possível distinguir alguma marcação especial para as letras em início de frase. - Cada desvio foi considerado como uma ocorrência, mesmo que dois desvios fossem relativos ao mesmo aspecto linguístico.
A grade de avaliação do desempenho linguístico considerou, portanto, três grandes grupos de competências, segundo os aspectos explicitados anteriormente : a ) domínio das convenções ortográficas : grafia de vogais e consoantes, uso de maiúsculas e minúsculas, emprego do hífen e acentuação gráfica ; b ) domínio dos procedimentos de estruturação textual do ponto de vista microestrutural : organização interna dos períodos, emprego de conectores para a articulação lógica entre os períodos e entre os parágrafos, emprego de marcas de referenciação lexical e pronominal ; utilização dos sinais de pontuação que contribuem para a organização lógica da frase ; c ) domínio das regras de caráter morfossintático estabelecidas como modelares do ponto de vista da modalidade escrita formal da norma-padrão da Língua Portuguesa : concordância nominal e verbal, regência nominal e verbal, colocação pronominal, flexão nominal e verbal, correlação entre tempos e modos verbais, ausência de marcas de oralidade.
A seleção vocabular adequada à modalidade escrita formal da Língua Portuguesa foi incorporada a essa última competência, tendo em vista a intersecção entre as duas do ponto de vista das exigências do registro formal da modalidade escrita da norma-padrão.
Apreciação geral do desempenho dos estudantes : A correção revelou desempenhos distintos dos participantes, marcados pela falta de repertório cultural da maioria dos participantes : - a questão 1 propiciou a oportunidade de maior desenvolvimento do tema solicitado, gerando consequências na estruturação textual, já que os textos foram mais longos, mais elaborados e fluentes.
Entretanto, em virtude da diversidade de abordagens dos textos motivadores ( a falta de penicilina, a recusa no uso da camisinha pelos homens, a maior frequência das mulheres na realização de exames ginecológicos, a transmissão das DST aos bebês durante a gestação ), observaram -se respostas que se configuraram como verdadeiras paráfrases, com comprometimento do caráter \textbf{autoral} desejável em uma questão discursiva de formação geral.
Alguns participantes selecionaram um trecho de cada texto motivador para compor suas respostas, outros copiaram trechos inteiros. - a questão 2 teve um comportamento distinto, devido ao caráter polêmico que o caracteriza.
Assim, os textos são mais curtos, com vocabulário repetitivo e preso aos textos motivadores.
Além disso, observou -se grande ocorrência de protesto em função do tema, ora por motivos religiosos ora por preconceito de gênero.
Quanto aos aspectos linguísticos analisados durante esta avaliação, os resultados observados estão descritos nos parágrafos seguintes.
Aspectos ortográficos : O desempenho dos participantes revelou uma diferença muito grande nos dois aspectos analisados nesta competência : baixo índice de desvios da grafia padrão e grande índice de desvios de acentuação.
Em vários casos, ocorre ausência completa de acentuação gráfica.
Os resultados revelam que a tendência dominante entre os universitários brasileiros é a eliminação da acentuação gráfica, talvez motivada pelos hábitos relacionados às redes sociais e pela ausência de esclarecimento dos meios de comunicação, das autoridades e das escolas sobre as decisões do Acordo Ortográfico da Língua Portuguesa de 1990.
Os casos mais sistemáticos de eliminação do acento indicador da sílaba tônica são :  palavras proparoxítonas ( “ sifilis ”, “ proximos ”, “ politicas ”, “ publicas ”, “ transgeneros ” ) ;  palavras paroxítonas terminadas em ditongo crescente ( “ necessario ”, “ noticiarios ”, “ individuo ”, “ dependencia ”, “ varios ”, “ propria ”, “ transmissiveis ”, “ ocorrencia ” ) ;  palavras paroxítonas com hiato ( “ saude ” )  palavras oxítonas ( “ ninguem ”, “ esta ”, “ ate ”, “ tambem ”, “ prevencao ” ).
Por outro lado, destaca -se o uso indevido do acento gráfico em palavras como “ genêro ”, “ melâncolia ”, “ prevênção ”, “ intervênção ”, por exemplo.
Há, também, uma tendência a acentuar paroxítonos como se fossem oxítonos : “ transmissivéis ”.
Quanto ao domínio das convenções relativas à grafia das palavras, observam -se desvios recorrentes, como a hipercorreção pela escolha de “ e ” no lugar de “ i ”, por influência de hábitos da oralidade ( “ descriminação ” no lugar de “ discriminação ”, “ entervenção ” no lugar de “ intervenção ” ) ; desvios de grafia relacionados à variação diastrática podem ser observados em “ estrupo ”, “ subjulgadas ”, “ precoseito ”, “ soubre ”, “ vecendo ”, “ indesencia ”, “ apolojia ”, “ fulga ”, “ dereitos ”, “ sifelis ”.
Observam -se, também, casos de inadequação no uso da maiúscula para destacar determinadas palavras-chave do texto, como “ Transgêneros ”, “ Brasileiros ”, “ Homens ”.
Destaque -se, também, o grande número de participantes que grafam os textos inteiramente em caixa alta.
Vale observar, também, que, ao contrário do que se esperava, não apareceram abreviaturas próprias do “ internetês ”, relacionadas ao uso de redes sociais e e-mails.
Aspectos textuais : Esta competência é a que se revela como a mais problemática entre os participantes, porque são muitos os problemas observados, desvios acumulados durante toda a formação do estudante e que não se resolvem com um estudo autodidata, como acontece com regras ortográficas ou morfossintáticas.
São eles : sequência justaposta de ideias sem encaixamentos sintáticos ; redução drástica de estruturas subordinadas, ao lado do aumento na frequência de estruturas coordenadas e absolutas ; redução no uso de conectores para expressar relações lógicas essenciais à construção do texto, substituídas pela exigência de inferência por parte do interlocutor para suprir a sua ausência ; emprego equivocado de operadores que não estabelecem relações lógicas coerentes entre ideias do texto ; emprego inadequado do pronome relativo ( com omissão da preposição ou a utilização de pronome inadequado, como “ onde ” ) ; repetição exaustiva de termos sem a utilização de procedimentos mais sofisticados de substituição ( hiperonímias, hiponímias, nominalizações, expressões metafóricas ) ; frases fragmentadas que comprometem a estrutura lógico-gramatical ; frases formadas apenas por oração subordinada, sem oração principal.
Um importante aspecto a destacar é o baixo desempenho de uma parte dos participantes em relação à estrutura formal do texto produzido, o que é preocupante ao se levar em conta que são graduandos em fase final de formação.
São frequentes os casos de desvios de estruturação frasal, com uso inadequado ou ausência de conectivos entre parágrafos e entre frases.
Em uma parte dos textos, falta textualidade e domínio do registro padrão da língua.
Na verdade, observam -se relações linguísticas quase agramaticais, como as estabelecidas pela sequência de gerúndios sem o apoio de um ponto de partida para a organização das informações gramaticais e semânticas.
Observou -se que uma parte dos participantes não distribuiu as ideias em parágrafos, talvez devido ao pequeno número de linhas disponibilizadas para a resposta da questão ou, quem sabe, pela suposição de que não seria necessária essa divisão por não se tratar de um texto no modelo de uma redação dissertativo-argumentativa, como solicitado nos vestibulares.
Em função do encaminhamento dos enunciados das duas questões, que solicitaram uma análise do problema e encaminhamentos de políticas públicas, houve uma grande tendência, também, de construção de dois parágrafos desconexos, sem utilização de elementos coesivos adequados à progressão textual.
Quanto à utilização dos mecanismos de referenciação, deve -se destacar a ocorrência de repetições de palavras ou expressões sem a utilização de termos sinônimos ou pronomes, como seria adequado.
Quanto à utilização dos sinais de pontuação, observou -se uma precariedade nos textos analisados.
É muito frequente a ocorrência de parágrafos sem marca interna de pontuação para separar os períodos.
Vale observar que não foi penalizada a ausência de vírgula para destacar locuções ou adjuntos adverbiais de pequena extensão deslocados de posição na frase, por ser um uso opcional.
São os seguintes os tipos de problemas encontrados : a ) vírgula : utilização de vírgula para separar o sujeito e o predicado ; ocorrência de apenas uma das vírgulas para separar uma palavra, uma expressão ou uma oração encaixada ; uso de vírgula no lugar do ponto para separar ideias que constituem períodos distintos ; ausência de vírgula para separar elementos de uma enumeração ; ausência de vírgula para separar oração adjetiva explicativa ou utilização inadequada para separar oração adjetiva restritiva ; b ) ponto e vírgula : utilização do ponto e vírgula no lugar de vírgula ; c ) ponto final : ausência de ponto final para separar períodos.
Aspectos morfossintáticos e vocabulares : Em relação à regência, o desvio mais frequente é a falta do sinal indicativo da crase – isso revela que o usuário não tem consciência de que, sob a forma do termo “ a ”, existe a presença de uma contração entre a preposição “ a ” ( exigida pela regência do termo anterior ) e o artigo definido “ a ”.
Um desvio de regência significativo, nos últimos anos, é a utilização inadequada de uma preposição ou sua ausência após o verbo ou o nome ( substantivo ou adjetivo ).
Outro problema relacionado à regência verbal e à nominal, encontrado frequentemente nas questões, foi a ausência de preposição antes de pronome relativo, processo generalizado na modalidade oral da língua, em situações de registro informal.
Apesar da possibilidade de que essa alteração de regência se generalize no padrão escrito da Língua Portuguesa, como já está ocorrendo até em textos jornalísticos, o não emprego da preposição foi considerado inadequado neste processo de avaliação.
A concordância verbal e a concordância nominal apresentam alguns desvios muito frequentes.
Quanto à concordância de número, observou -se ausência de marca ( com sujeito anteposto ou posposto ) ou uso indevido ( uso inadequado da marca de plural comandado pelo núcleo plural da locução adjetiva, apesar de o substantivo que funciona como núcleo do sintagma nominal estar no singular ).
Uma ocorrência que se destacou foi a ausência de acento circunflexo na forma plural do presente do indicativo do verbo “ ter ”, que foi considerada como um desvio na concordância verbal e não na acentuação gráfica.
Quanto à concordância de gênero, vários casos foram observados, normalmente no âmbito de sintagmas nominais longos, em que o adjetivo está afastado do substantivo.
Deve -se destacar o aparecimento da marca de plural em verbos ou adjetivos comandados por adjunto adnominal plural, apesar de serem relacionados a núcleos substantivos no singular, evidenciando um processo de hipercorreção ( exemplo : “ O nome social dos trangêneros podem ser estabelecidos por uma legislação específica ” ).
Quanto à questão da colocação pronominal, foram poucos os desvios observados.
Concluiu -se que, no registro escrito formal, a maioria dos participantes já incorporou regras como a não introdução da frase por um pronome oblíquo e a próclise na presença de um termo atrator.
Não se adotou, entretanto, o padrão excessivamente formal descrito pelas gramáticas normativas em relação à posição do pronome oblíquo em locuções verbais, já que esse uso está muito distante da prática cotidiana, até em textos mais formais.
Quanto aos aspectos vocabulares, alguns tipos de inadequação foram observados : expressões da oralidade ; seleção vocabular incompatível com o contexto, gerando falta de inteligibilidade ; falta de domínio de vocabulário mais abstrato e de maior complexidade, essencial ao desenvolvimento do texto de base dissertativa.
O principal aspecto observado foi a excessiva repetição de certas palavras, revelando limitação de repertório vocabular.
O termo “ pessoa ”, por exemplo, chega a ser repetido até 7 ou 8 vezes em um mesmo texto.
Várias marcas de oralidade foram identificadas, embora não com alta frequência : uso do pronome relativo “ onde ” como relativo universal, falta de artigo definido antes de substantivo, reduções como “ tá ”, “ pra ”, “ pro ”, “ prum ”, expressões informais, eliminação de preposições.
Em função do tema solicitado na questão 1, a sigla DST foi grafada de diferentes maneiras : DST´s, DSTs ou DSTS.
Além disso, houve flutuação de gênero gramatical nos sintagmas que envolviam a sigla : “ os diferentes DSTs ” ou “ as diferentes DSTs ”.
Em função do tema solicitado na questão 2, houve flutuação de gênero gramatical nos sintagmas que envolviam a sua palavra-chave : “ os trans ” ou “ as trans ” ; “ os transgêneros ” ou “ as transgêneras ”.
Todas essas formas foram aceitas, por considerarmos que sua utilização é muito recente e os participantes ainda não sistematizaram essas formas em seu uso da língua. 6.3.2 Componente de Conhecimento Específico Na parte da prova relativa às questões discursivas do Componente de Conhecimento Específico ( Tabela 6.17 ), observa -se que a Média foi mais baixa do que para as questões discursivas do Componente de Formação Geral.
Enquanto no Componente de Formação Geral a Média para estudantes de Arquitetura e Urbanismo de todo o Brasil foi 52,2, na parte de Conhecimento Específico a Média foi 36,5.
A maior Média deste componente foi obtida pelos estudantes da região Sul ( 37,5 ), e a menor, pelos da região Norte ( 30,3 ).
Quanto à variabilidade das notas, o Desvio padrão de todo o Brasil foi 19,7.
O maior Desvio padrão foi encontrado na região Nordeste ( 20,8 ), e o menor, nas regiões Norte e Sul ( 19,1 ).
A maior nota Máxima, 93,3, foi obtida na região Sudeste.
Nas demais regiões a nota Máxima foi : 90,0 nas regiões Norte, Nordeste e Centro-Oeste e 86,7 na região Sul.
A nota Mínima ( 0,0 ) foi obtida por pelo menos um aluno em todas as regiões do Brasil.
A Mediana do Brasil como um todo foi 36,7.
A maior Mediana foi obtida nas regiões Nordeste e Sul ( 40,0 ), e a menor, na região Norte ( 30,0 ).
O Gráfico 6.10 representa a distribuição das notas nas questões discursivas no Componente de Conhecimento Específico.
A moda desta distribuição ocorre no intervalo [ 30 ; 40 ], com em torno de 18 \% do total de participantes.
A frequência desse intervalo é seguida de perto pela do intervalo ( 40 ; 50 ].
Destaca -se, também, o intervalo [ 0 ; 10 ], que inclui, além das notas até 10, a frequencia de estudantes que deixaram a resposta às três questões em branco.
A análise de cada uma destas questões será feita a seguir.
O Gráfico 6.10 representa a distribuição das notas nas questões discursivas no Componente de Conhecimento Específico.
A moda desta distribuição ocorre no intervalo [ 30 ; 40 ], com em torno de 18 \% do total de participantes.
A frequência desse intervalo é seguida de perto pela do intervalo ( 40 ; 50 ].
Destaca -se, também, o intervalo [ 0 ; 10 ], que inclui, além das notas até 10, a frequencia de estudantes que deixaram a resposta às três questões em branco.
A análise de cada uma destas questões será feita a seguir. 6.3.2.1 Análise da Questão Discursiva 3 do Componente de Conhecimento Específico Na questão 3, cujos resultados aferidos encontram -se descritos na Tabela 6.18, a Média dos estudantes de todo o Brasil foi 46,0, a questão de melhor desempenho dentre as três discursivas de Conhecimentos Específicos.
A menor Média nessa questão foi obtida pelos alunos da região Norte ( 37,1 ), enquanto a maior Média foi obtida na região Sudeste ( 47,0 ).
Quanto à variabilidade das notas, o Desvio padrão de todo o Brasil foi 28,7.
O maior Desvio padrão foi obtido na região Centro-Oeste ( 29,7 ), enquanto o menor foi obtido na região Sul ( 28,0 ).
A nota Máxima, 100,0 pontos, foi alcançada por, pelo menos, um aluno de todas as regiões.
A Mediana do Brasil como um todo foi 50,0, a mesma para quatro das Grandes Regiões.
A exceção foi a Mediana da região Norte, igual a 40,0.
A nota Mínima ( 0,0 ) foi a mesma em todas as regiões do Brasil.
O Gráfico 6.11 mostra a distribuição das notas na questão discursiva 3, do Componente de Conhecimento Específico da área de Arquitetura e Urbanismo.
Pouco menos do que 10 \% dos estudantes deixaram a questão em branco. maior frequência, moda da distribuição, encontra -se no intervalo ( 50 ; 60 ], com cerca de 18 \% dos participantes nesse intervalo.
O Gráfico 6.11 mostra a distribuição das notas na questão discursiva 3, do Componente de Conhecimento Específico da área de Arquitetura e Urbanismo.
Pouco menos do que 10 \% dos estudantes deixaram a questão em branco. maior frequência, moda da distribuição, encontra -se no intervalo ( 50 ; 60 ], com cerca de 18 \% dos participantes nesse intervalo. 6.3.2.2 Comentários sobre as respostas à Questão Discursiva 3 Com relação ao enunciado da Questão 3, três aspectos merecem reflexão.
O primeiro aspecto diz respeito ao conceito de"habitação social", objeto central da questão.
Talvez a apresentação da definição desse conceito no enunciado evitasse \textbf{repostas} inconsistentes.
Embora grande parte dos graduandos da Área considere esse termo relativo a unidades habitacionais produzidas com financiamento ou subsídio governamental e destinadas a moradores de rendas baixas, registraram -se respostas que não consideraram o termo dessa forma.
Dentre esses, uma pequena parte dos formandos interpretou o conceito como unidades de moradia destinadas à população em geral, independentemente da faixa de renda, ou até mesmo associando os tipos citados à população de rendas mais elevadas.
O segundo aspecto, que levou a várias interpretações diferentes da esperada, foi a descrição dos modelos tipológicos de habitação de interesse social.
Um deles, segundo o enunciado, “... caracteriza -se por conjuntos de habitação multifamiliar verticalizada, que ocupa lotes de grandes dimensões, cercados por muros ”.
Nesse caso, parte das respostas se ateve ao fato de ser habitação vertical, considerando como negativa a presença de edifícios altos e de muros, correlacionando -se com “ condomínios fechados ” de média e alta renda, que não era o objeto da questão.
O outro modelo foi descrito como “ conjuntos de habitação unifamiliar isolada no lote ”, e detectou -se um equívoco de interpretação do termo"isolada".
Em algumas respostas, o termo foi interpretado como um conjunto “ isolado ” e não como uma edificação afastada das divisas do terreno onde se insere, e que é um modelo frequentemente utilizado para moradias unifamiliares.
O terceiro aspecto diz respeito ao comando da questão.
Foi solicitada a redação de um texto com a descrição de duas consequências da aplicação dos modelos, “ relacionadas à qualidade do espaço urbano ou à infraestrutura urbana ”, e não havia indicação se as consequências deveriam ser “ negativas ” ou “ desvantagens ”, como a maioria das apresentadas no padrão de respostas.
Algumas respostas indicaram consequências positivas, pois, em muitas situações, os formandos assim interpretam os conjuntos habitacionais, lembrando que alguns desses formandos residem em moradias desse tipo.
Pelo exposto, a questão possibilitou interpretações diferentes da esperada pelo padrão de resposta.
Observou -se que as respostas se subdividiram em grupos distintos, provavelmente por se condicionarem ao perfil sociocultural do(a) formando(a), que, muitas vezes, habita uma unidade de habitação de interesse social, ou ao enfoque dado ao tipo de habitação, vertical ou horizontal, ou ainda à presença de edifícios altos e muros, sem citar a condição de interesse social.
O espaço para apresentação da resposta pareceu insuficiente em muitos casos.
A maioria dos estudantes preencheu todo o espaço fornecido ou até utilizou margens do espaço disponível, provavelmente por se tratar de um tema conhecido, com vertentes diversas e conhecimento amplo por parte dos(as) formandos(as) na Área.
O conteúdo abordado é adequado ao formando da Área de Arquitetura e Urbanismo, e a dificuldade da questão era baixa, especialmente em comparação às outras questões.
Em relação à diversidade de respostas encontradas, como a questão tem caráter qualitativo e é de base geral, sem se ater em contextos específicos, o grau de flexibilidade foi elevado.
Foi grande a incidência de respostas previstas no padrão, mas outras respostas não previstas, algumas corretas, foram também encontradas.
Dentre essas, conforme comentado no item anterior, alguns destacaram consequências positivas.
Outros compararam os tipos horizontais ou verticais, sem se ater ao perfil de interesse social da habitação, ou focaram a existência de"muros"ou"unidades isoladas", também sem se referir ao perfil de interesse social da habitação.
Observaram -se respostas incompletas, por desconhecimento do tema, e ainda algumas respostas que repetiram conteúdo do enunciado, apenas parafraseando -o.
Cabe destacar que, de forma geral, houve a relação entre consequências e explicações possíveis, conforme padrão de resposta, embora a relação com os modelos de habitação de interesse social não tenha sido tão frequente.
A linguagem adotada pelos ( as ) formandos(as), de forma geral, foi adequada para a Área de Arquitetura e Urbanismo, demonstrando conhecimento de aspectos relativos às consequências da implantação intensiva de conjuntos de unidades multifamiliares e unifamiliares de habitação de interesse social.
Faltaram respostas que considerassem o grande diferencial entre os modelos : ser fruto de processos de desmembramento de glebas sob o formato de condomínio ( unidades multifamiliares ) ou ser resultado de projetos de loteamento tradicional ( unidades unifamiliares ).
É possível que a falta de respostas que considerassem o diferencial entre os modelos se deva ao fato de que a questão fundiária, fundamental na forma de acesso à terra urbana, não estar sendo abordada adequadamente nos currículos das escolas.
As abordagens mais frequentemente encontradas para responder incluíam : a segregação socioespacial como consequência da aplicação dos modelos de habitação de interesse social ; a localização periférica dos conjuntos habitacionais tanto uni como multifamiliares ; a maior demanda por deslocamentos e a precariedade dos sistemas de transporte público ; a carência de infraestrutura física e social ; o estigma dos moradores ; a deficiência das políticas públicas, de forma geral.
Como no padrão de resposta foi solicitado um texto sobre duas consequências relacionadas à qualidade do espaço urbano ou à infraestrutura urbana, foi grande a diversidade de respostas encontrada.
Os principais equívocos se ativeram ao desconhecimento do termo"habitação de interesse social", como associado ao perfil de renda baixa de seus moradores ou ao fato de ser financiada ou subsidiada pelo setor público e ao desconhecimento dos tipos de habitação de interesse social descritos : vertical ( conjuntos de blocos de até 4 pavimentos em condomínios ) e horizontal ( unidades isoladas no lote e que se referem a casas térreas, unifamiliares em loteamentos abertos ).
Na pontuação, as avaliações que apresentaram nível mediano a alto se relacionaram às respostas que se referiram mais diretamente à habitação de interesse social, conforme indicado no enunciado, e que descreveram : duas ou mais consequências com suas possíveis explicações, referências corretas aos tipos de habitação ( vertical e horizontal ) e relação com aspectos de infraestrutura ou espaço urbano ; ou uma consequência com possível explicação e uma consequência sem explicação, referências sobres os tipos de habitação ( vertical e horizontal ) e relação superficial com aspectos de infraestrutura ou espaço urbano.
As respostas que receberam pontuações menores se relacionaram à descrição da habitação, independentemente da destinação à população de renda baixa, ou que não interpretaram corretamente os tipos de habitação ( vertical e horizontal ) ou que não fizeram relação com aspectos de infraestrutura ou espaço urbano.
A Questão 3 descreve um aspecto fundamental para o ambiente urbano das cidades brasileiras, relativo à produção de habitação de interesse social no Brasil, sendo recorrente nas situações em que o(a) arquiteto(a) e urbanista poderá se deparar na prática de sua profissão, como gestor(a) público(a), profissional autônomo(a) contratado(a) ou empresário(a ).
É possível que as deficiências na formação tenham relação com a conceituação sobre habitação de interesse social e ao histórico desse tema no Brasil, e de certa forma, indicando que se faz necessária, uma discussão mais aprofundada sobre tipos morfológicos de unidades habitacionais de interesse social, no que concerne às distintas modalidades de empreendimentos por desmembramento de glebas parceladas, sob a modalidade de condomínio ou por parcelamento, sob a forma de loteamento.
Essa diferenciação está na base de formação tipológica dos conjuntos habitacionais. 6.3.2.3 Análise da Questão Discursiva 4 do Componente de Conhecimento Específico A Tabela 6.19 contém as informações relativas à questão 4 do conjunto de questões do Componente de Conhecimento Específico.
O desempenho dos estudantes nessa questão foi inferior ao das questões 3 e um pouco superior ao da questão 5.
A Média geral do Brasil foi 32,1, sendo a menor Média registrada na região Norte ( 26,8 ), e a maior, na região Nordeste ( 33,2 ).
A nota Máxima ( 100,0 ) foi atingida por, pelo menos, um concluinte de todas as regiões.
A Mediana em todo o Brasil foi 30,0, o mesmo valor foi obtido em quatro Grandes Regiões, sendo a Mediana da região Nordeste a única diferente ( 40,0 ).
A nota Mínima foi zero em todas as regiões.
O Gráfico 6.12 representa a distribuição de notas da questão discursiva 4, do Componente de Conhecimento Específico.
Essa distribuição tem moda na classe de estudantes que deixaram a resposta à questão em branco, com frequência pouco superior a 15 \%.
Desconsiderando os que deixaram a resposta em branco, a maior frequência ocorre no intervalo ( 30 : 40 ], seguido de perto pelo intervalo ( 40 ; 50 ], todos os dois com frequência pouco inferior a 15 \%. 6.3.2.4 Comentários sobre as respostas à Questão Discursiva 4 A questão 4 abordou uma questão fundamental para o ensino de projeto urbano e arquitetônico no momento atual no Brasil, intervenções urbanas participativas em projetos que tratam de habitação social e ações de planejamento de forma geral.
A relevância também reside no valor do patrimônio local.
Trouxe à tona um caso específico : a Vila Itororó, cujo tombamento municipal é de 2002, o estadual de 2005 e que, em 2006, foi decretada como área de utilidade pública.
Além do aspecto de como a população pode atuar em prol da valorização do patrimônio com ações culturais, destacam -se aspectos sobre a ocupação e a função social da propriedade em paralelo ao processo participativo de gestão.
Os conhecimentos exigidos pelo enunciado estavam de acordo com as Diretrizes Curriculares da Área no que diz respeito às disciplinas de Projeto de Arquitetura e aos procedimentos de ensino adotados nestas disciplinas.
O nível de profundidade esperado pelo padrão de resposta foi alto para o que se espera de um(a) graduando(a) em Arquitetura e Urbanismo, por três motivos : - As diferentes configurações dos cursos por todo o Brasil, que nem sempre abordam de modo explícito as implicações que determinados processos projetuais podem provocar ; - A falta de incentivo à descrição e/ou verbalização dos procedimentos projetuais durante as disciplinas de Projeto de Arquitetura ; - A falta de adequação da linguagem adotada pelos alunos em relação aos termos específicos da Área.
Importante ressaltar que a valorização do patrimônio cultural também é uma das atribuições do Arquiteto e Urbanista, estabelecida pelo Conselho de Arquitetura e Urbanismo ( CAU ).
A questão foi considerada pela banca mais difícil em comparação com as questões 3 e 5, por focalizar projeto e contexto específico, além de possibilitar diferentes interpretações devido à inserção de imagens referentes aos processos de projeto discutidos na questão.
Embora a linguagem fosse simples, as informações e os dados e imagens, não muito claras, apresentados no enunciado para a solução do problema, não se mostraram suficientes para esclarecer ao aluno o que realmente era solicitado, principalmente por abordar um caso específico de uma grande cidade brasileira ( São Paulo ).
A realidade descrita é mais comumente debatida em escolas de Arquitetura e Urbanismo localizadas em grandes cidades e capitais do Sudeste, principalmente em São Paulo, onde se situa o estudo de caso apresentado na questão.
O possível desconhecimento do sítio por estudantes de outras regiões do país, talvez tenha contribuído para interpretações diferentes da esperada pelo padrão de resposta.
Dois tipos de respostas divergentes da esperada foram mais recorrentes : - Aquelas que distinguiram somente diferenças estéticas e formais entre os dois conjuntos de imagens ; - Aquelas que fizeram a ligação direta entre a derrocada dos projetos e das obras públicas e os casos de corrupção e desvios de verbas públicas atualmente em evidência no Brasil.
A dificuldade de expressão dos conhecimentos, por vezes evidente, demonstrou a deficiência na formação do(a) estudante de Arquitetura e Urbanismo no que diz respeito ao trabalho com a interpretação e comunicação textual.
Foi observado que as respostas se subdividiram em grupos distintos por se condicionarem : ( i ) à característica e uso das edificações do conjunto residencial abordado, com foco na atividade cultural, que induziu a respostas sobre a arquitetura resultante, ilustrada nas imagens ; ( ii ) ao impacto no entorno físico e social, principalmente em relação à atividade como centro cultural, sendo que algumas das respostas demonstraram conhecimento do conjunto ; ( iii ) aos procedimentos metodológicos em si, sem citar a nomenclatura esperada.
É possível que os enganos cometidos em respostas que nem sequer demonstraram saber que existem diferenças entres os dois processos projetuais apresentados, em que pese a dificuldade de interpretação do enunciado e das imagens, correspondam às deficiências na formação dos alunos em relação ao conteúdo avaliado pela questão.
Em função das \textbf{observações} já mencionadas, houve um conjunto de interpretações que geraram respostas que não atenderam à definição esperada de cada metodologia ( “ compartilhada ” e “ centralizada ” ), mas que, em parte, identificaram as diferenças entre elas.
Foram citadas algumas implicações de cada método, e algumas contrariaram a expectativa de resposta que indicava uma metodologia ( “ compartilhada ” ) com melhores implicações para os moradores do que a outra ( “ centralizada ” ).
Em relação às diferenças de pontuação : - As respostas que foram consideradas insuficientes apresentaram basicamente a identificação de que havia dois tipos diferentes de projetos nos dois conjuntos de imagens ou somente identificaram que em um procedimento não havia a participação da população e no outro havia ; - Nas respostas consideradas medianas havia uma caracterização nominal dos dois diferentes processos projetuais e uma breve explicação sobre as implicações, não necessariamente sobre os dois ; - As respostas boas distinguiram e nomearam os dois processos e apresentaram explicações mais consistentes sobre as implicações de cada um deles.
Uma dificuldade revelada em algumas respostas estava associada ao fato de não ter sido informado o ano da construção original da Vila Itororó.
Como o texto mencionava o ano de 1974, como o ano da primeira proposta, alguns alunos tomaram esse ano como sendo o do projeto original.
A carência de descrição sobre o tipo de habitação ao qual a Vila Itororó se destinava originalmente, seu valor patrimonial local, o tempo em que o conjunto ficou abandonado e sobre quando foi ocupado irregularmente por moradores, possivelmente contribuiu para a dificuldade da elaboração da resposta.
Outro aspecto a ser apontado foi o uso das imagens como fonte de informação e base para resposta, uma vez que o foco da pergunta é o processo de projeto e não o resultado.
Ao pedir que os estudantes considerassem “ as imagens e as informações apresentadas ”, muitas respostas se pautaram na característica das edificações, organização espacial, estilo arquitetônico, dentre outros aspectos.
Além disso, as imagens foram apresentadas sem contextualização/legenda, apenas a fonte, sem mencionar os autores do projeto, o que é inadequado.
As legendas seriam fundamentais para o melhor entendimento da situação.
A omissão de informação sobre o projeto original, o abandono, os projetos de restauro e as ações participativas, o canteiro aberto talvez tenham induzido a respostas difusas, fora do eixo central da pergunta – o processo de projeto, para o qual as informações não fornecidas são relevantes.
Cabe ressaltar que em nenhuma das respostas foi identificado que o procedimento metodológico “ centralizado ” correspondeu a um projeto realizado em 1974, período de ditadura militar no país.
Essa reflexão poderia ser debatida por se tratar de período recente na história do país.
Outra \textbf{observação} diz respeito à demanda de se falar da questão patrimonial, considerando que técnicas retrospectivas estão nas Diretrizes Curriculares, tal como a valorização do patrimônio cultural, também, uma das atribuições do Arquiteto e Urbanista estabelecida pelo Conselho de Arquitetura e Urbanismo-CAU. 6.3.2.5 Análise da Questão Discursiva 5 do Componente de Conhecimento Específico A Tabela 6.20 contém as informações relativas à questão 5 do conjunto do Componente de Conhecimento Específico.
O desempenho dos estudantes de todo o Brasil nesta questão foi o mais baixo dentre as três questões discursivas desse componente.
A nota Média dos estudantes de todo o Brasil foi 31,5.
A maior Média foi registrada na região Sul ( 33,0 ), enquanto a menor Média foi registrada na região Norte ( 26,9 ).
Quanto à variabilidade das notas, o Desvio padrão dos alunos do Brasil, como um todo, foi 24,4.
Enquanto o maior desvio foi encontrado na região Nordeste ( 25,7 ), o menor foi encontrado na região Norte ( 23,7 ).
A Mediana para o Brasil e em todas as regiões foi 30,0.
Para o conjunto de alunos de Arquitetura e Urbanismo do Brasil, a nota Máxima foi 100,0, e a nota Mínima ( 0,0 ).
Os mesmos limites foram encontrados em todas as regiões.
O Gráfico 6.13 representa a distribuição das notas da questão discursiva 5 do Componente de Conhecimento Específico.
Essa distribuição, como a relativa à questão discursiva 4, também tem moda na categoria “ em branco ”.
Destaca -se que a frequência de notas zero somada a de estudantes que não responderam à questão \textbf{ultrapassa} 20 \% do total de participantes.
Além desse grupo modal, observa -se um máximo local no intervalo ( 20 ; 30 ], seguido de perto pelo intervalo ( 30 ; 40 ]. 6.3.2.6 Comentários sobre as respostas à Questão Discursiva 5 O enunciado da questão trazia uma adaptação do trecho da cartilha da Lei no 12.587 de 2012, que trata da Política Nacional de Mobilidade Urbana, problemática fundamental para o campo do planejamento urbano, envolvendo a acessibilidade urbana como política pública e direito constitucional.
Apesar de abordar tema atual e fundamental do urbanismo, dentre as três questões discursivas específicas, as respostas refletiram que essa foi a que gerou maiores dificuldades em relação ao conteúdo.
Observou -se que as respostas se relacionaram a quatro fatores mais recorrentes : ao acúmulo de carros particulares como sendo um dos problemas relacionados à mobilidade, como trânsito caótico ; ao meio-ambiente, devido aos impactos causados na qualidade do ar e na poluição sonora ; às deficiências do transporte coletivo, pincipalmente no tocante ao funcionamento de ônibus ; e à necessidade de intensificar e estruturar o uso de bicicletas e o caminhamento de pedestres.
Além dos aspectos que se referiram basicamente às questões de transporte, foram indicados outros relacionados ao espaço urbano, demonstrando, em algumas situações, a capacidade e maturidade de \textbf{observação} de questões integradas à mobilidade que transcendem elementos de transporte, público ou privado, motorizado ou não motorizado.
Foi frequente a perspectiva de que mobilidade urbana sustentável é uma política pública que por si resulta numa qualificação do espaço urbano.
Em algumas das respostas foram consideradas outras políticas que deveriam ser integradas ( saneamento, habitação, acesso à terra, etc. ) para se obter melhor qualidade de vida e mais inclusão social nas cidades.
Cabe destacar que alguns(as) estudantes responderam de forma bem interessante à questão, trazendo conceitos de integração de demais políticas sobre uso e ocupação do solo, arborização, implantação de parques e áreas verdes, desenho urbano, compacidade e densidade de tecido urbano, entre outras, como forma de responder à mobilidade e à acessibilidade.
Dentre as respostas diversas dadas pelos(as) formandos(as) podem ser citadas : a ) identificação de aspectos jurídicos sobre direito à cidade e à mobilidade, ou a argumentação sobre aspectos referentes a políticas públicas que envolvem direta ou indiretamente a mobilidade ; b ) aspectos relacionados a transportes em si, tanto públicos como privados ( trânsito, problemas relacionados aos carros particulares, vantagens de bicicletas, implantação de linhas de metrô e VLT, melhorias no transporte de ônibus, etc. ) ; c ) aspectos relacionados a demais elementos do espaço urbano que se relacionam direta ou indiretamente à mobilidade ( densidade habitacional, renda, investimentos financeiros, uso do solo, áreas verdes, parques, arborização de calçadas, etc. ).
A linguagem utilizada pelos(as) formandos(as) foi adequada com base nos termos específicos da Área de Arquitetura e Urbanismo, embora novamente deva -se ressaltar que parte das respostas continha basicamente textos do enunciado.
Dentre as respostas não previstas pelo padrão podem ser incluídas as que citaram aspectos relacionados a demais elementos do espaço urbano, tais como : questões financeiras e de renda da população ; densidade de ocupação ; uso funcionalmente diversificado do solo urbano ; implantação de áreas verdes e parques ; arborização de ruas e calçadas ; dispositivos de acessibilidade e desenho universal.
Esses aspectos contribuem para uma maior diversidade de usos e funções do espaço urbano e para maior conforto aos pedestres, devido ao melhor usufruto do ambiente urbano levando à intensificação de circulação de pedestres, ciclistas e pessoas com deficiência. 6.3.3 Considerações Finais A coordenação da correção das três questões discursivas da Área de Arquitetura e Urbanismo foi bastante positiva para reflexão sobre os conteúdos temáticos abordados nas questões.
A principal contribuição foi conhecer, através das respostas, as realidades distintas apresentadas pelos ( as ) formandos ( as ) na Área que diferem em forma, mas convergem em conteúdo.
Foi possível verificar que conceitos associados a modelos de habitação social, a processos participativos, como método de projeto, à mobilidade urbana e acessibilidade, estão presentes nos currículos dos cursos.
No entanto observamos que esses conceitos necessitam maior aprofundamento e integração aos conteúdos mínimos curriculares, a fim de melhor instrumentar os ( as ) estudantes em assuntos de vital importância para nossas cidades.
Em relação às questões e a seus enunciados, cabe ressaltar que as questões 3 e 5, por não se relacionarem a um caso específico, apresentaram melhor condição de resposta, embora tenham sido detectadas algumas dificuldades nos seus enunciados.
A questão 4, por sua vez, trouxe como fator \textbf{complicador} o fato de abordar um caso específico da cidade de São Paulo.
Cabem \textbf{observações} sobre a característica das questões dedicadas predominantemente a temáticas urbanas, recomendando -se em próximos exames que as questões discursivas sejam mais distribuídas entre temáticas relativas a projeto de edificações e temáticas relativas a projetos urbanos e paisagísticos, pois estão dentro do rol de atribuições, segundo o Conselho que regulamenta a profissão - Conselho de Arquitetura e Urbanismo.

% matched lemmas: advogado, assento, autoral, cartório, casal, complicador, observação, repor, ultrapassar
\end{textsample}
